\chapter*{Tóm tắt}
\addcontentsline{toc}{chapter}{Tóm tắt}

Báo cáo tìm hiểu bài báo của Kassem (2022) về đánh giá quản lý rủi ro trong dự án xây dựng dầu khí bằng mô hình phương trình cấu trúc bình phương tối thiểu riêng phần (PLS-SEM). Bài báo dùng dữ liệu khảo sát 314 người ở Yemen, xếp hạng 51 yếu tố rủi ro bằng chỉ số tầm quan trọng tương đối (RII) và ma trận xác suất -- tác động, rồi xây dựng mô hình hai bậc đo ảnh hưởng của rủi ro nội bộ và rủi ro bên ngoài tới thành công dự án.

Báo cáo trình bày cơ sở lý thuyết về quản lý rủi ro theo PMBOK và ISO~31000, cùng công thức và cách hoạt động của từng phương pháp. Vì dữ liệu gốc không được công bố, báo cáo tái hiện thực nghiệm theo hai hướng: kiểm tra tính nhất quán của số liệu đã công bố, và chạy lại toàn bộ quy trình trên một bộ dữ liệu mô phỏng có cấu trúc giống bài báo, với thuật toán PLS-SEM tự cài đặt bằng Python. Kết quả cho thấy phần tính toán cơ bản và kết luận chính của bài báo (rủi ro bên ngoài mạnh hơn rủi ro nội bộ) tái hiện được. Tuy nhiên, giá trị phân biệt không đạt ở một số nhóm rủi ro, và hệ số của 13 giả thuyết chi tiết không tái hiện được. Hai ví dụ thực tế, gồm dự án Nhiệt điện Thái Bình 2 và dự án Healthcare.gov, cùng phần minh họa bằng MS Project, Excel và mô phỏng Monte Carlo, cho thấy khung rủi ro của bài báo dùng lại được, nhưng mức độ quan trọng của từng nhóm rủi ro phụ thuộc mạnh vào bối cảnh.

\textbf{Từ khóa:} quản lý rủi ro, dự án xây dựng dầu khí, RII, ma trận xác suất -- tác động, PLS-SEM, HTMT, Monte Carlo.

\clearpage
